\documentclass[11pt,a4paper]{article}
\usepackage[utf8]{inputenc}
\usepackage[T1]{fontenc}
\usepackage{amsmath,amsfonts,amssymb}
\usepackage{graphicx}
\usepackage{hyperref}
\usepackage{listings}
\usepackage{xcolor}
\usepackage{booktabs}
\usepackage{float}
\usepackage{geometry}
\usepackage{algorithm}
\usepackage{algpseudocode}
\geometry{margin=1in}

\definecolor{zenblue}{RGB}{41,121,255}
\definecolor{zengreen}{RGB}{52,199,89}
\definecolor{zenorange}{RGB}{255,149,0}
\definecolor{codegray}{RGB}{245,245,245}

\hypersetup{
    colorlinks=true,
    linkcolor=zenblue,
    urlcolor=zenblue,
    citecolor=zenblue
}

\lstset{
    backgroundcolor=\color{codegray},
    basicstyle=\ttfamily\small,
    breaklines=true,
    captionpos=b,
    frame=single,
    numbers=left,
    numberstyle=\tiny\color{gray}
}

\title{
    \vspace{-2cm}
    \Large \textbf{Zen AI Model Family} \\
    \vspace{0.5cm}
    \Huge \textbf{Zen5: Flash} \\
    \vspace{0.3cm}
    \large Flash-Tier Frontier Inference, Behavioral Alignment, \\
    and Multimodal Routing on Personal Hardware \\
    \vspace{0.5cm}
    \normalsize Technical Whitepaper v2026.05
}

\author{
    Antje Worring\thanks{antje@hanzo.ai} \\
    Zach Kelling\thanks{zach@lux.network} \\
    \texttt{research@hanzo.ai} \\
    \\
    Zoo Labs Foundation \\
    \texttt{foundation@zoolabs.org}
}

\date{May 2026}

\begin{document}

\maketitle

\begin{abstract}
\textbf{Zen5} is not an independently trained model and not a native ``fifth generation'' of a Zen
foundation-model lineage. The Zen5 \emph{model} is a downstream repackaging of DeepSeek's open-weight
\textbf{DeepSeek V4 Flash} \cite{deepseekv4flash} (MIT) distributed as quantized GGUF; there is no
separate ``zen-5'' base model with its own weights, and any \texttt{zenlm/zen-5} base repository is a
placeholder---the only weights that exist are the DeepSeek V4 Flash GGUF. What we contribute is
packaging, behavioral post-training, and a multimodal routing layer, delivered as three coordinated
artifacts. First, the \emph{Zen5 inference engine}, a fork of antirez's \texttt{ds4} runtime
\cite{antirez-ds4} (MIT), serves the 284B-parameter DeepSeek V4 Flash mixture-of-experts model at
\texttt{IQ2\_XXS}/\texttt{Q2\_K} asymmetric quantization on a single 128 GB Apple Silicon machine. The engine reaches \textbf{27.30 tok/s
prefill and 8.93--26.68 tok/s generation} depending on prompt length, with a 1M-token context ceiling
and an on-disk SHA1-keyed KV cache that survives session and process restart. Second, a
\emph{behavioral post-training pipeline} consisting of (i) single-direction activation ablation against
the refusal subspace, (ii) identity supervised fine-tuning, and (iii) supervised fine-tuning on a curated
\emph{Zen Agentic Dataset} of tool-call traces produces a Zen5 persona that calls tools reliably and
does not over-refuse on legitimate dual-use prompts. Third, a \emph{multimodal routing} layer extends the
existing hanzoai/zen dynamic-router architecture so that Zen5 acts as the reasoning core while sibling
experts --- Zen-VL, Zen-3D, Zen-Foley, Zen-Musician, Zen-Video, and the Jin diffusion-transformer family ---
handle vision, depth, audio, music, video, and image generation, all colocated inside one
\texttt{hanzo-engine} process. As a forward research \emph{proposal} (not a shipped or evaluated
component) we describe \textbf{MoDE} (Mixture of Diverse Experts), a training-time diversity regularizer
for a diffusion MoE; the name refers to this proposed regularizer and is unrelated to any base-model
architecture. Zen5 ships as alpha-quality packaging code under MIT, with attribution to DeepSeek for the
V4 Flash weights (MIT) \cite{deepseekv4flash} and to \texttt{antirez/ds4} \cite{antirez-ds4} and
\texttt{ggml-org/llama.cpp} \cite{llamacpp} for the engine lineage.
\end{abstract}

\tableofcontents
\newpage

\section{Introduction}

The release of DeepSeek V4 Flash \cite{deepseekv4flash} and the corresponding \texttt{antirez/ds4}
inference engine \cite{antirez-ds4} established a new operating point for local AI: 284 billion
parameters running on a single laptop with 128 GB of unified memory, at usable wall-clock generation
speed, with a 1 million-token context window. This places frontier-tier reasoning on personal hardware
for the first time. As shipped, however, the system is intentionally narrow: it runs one specific GGUF
file, with one model identity, exposes one process talking one HTTP API on one port, and is tuned for
the persona and refusal patterns of the DeepSeek upstream. Zen5 closes this gap in three directions.

\paragraph{Engine.} We fork \texttt{antirez/ds4} as the Zen5 engine, preserving its sources, attribution,
and MIT license unchanged. Our deltas live in branding, distribution metadata, and a tracked set of
forthcoming features (Section~\ref{sec:engine}) that we will land upstream where possible and carry as
patches where they are Zen-specific.

\paragraph{Behavior.} The base model arrives with a generic DeepSeek persona and over-refuses on tasks
the Zen agent stack must perform: dual-use security analysis, edits in domain-restricted directories,
and faithful protocol-level reasoning during tool invocations. Zen5 applies a three-stage post-training
pipeline (Section~\ref{sec:posttrain}) that resolves identity, reduces refusal along the single-direction
subspace identified by Arditi et al.\ \cite{arditi2024refusal}, and improves tool-call fidelity through
supervised fine-tuning on the \emph{Zen Agentic Dataset}.

\paragraph{Modality.} A frontier reasoning model is necessary but not sufficient. Generation in 3D,
video, audio, music, and image is handled by specialized siblings within the Zen family --- Zen-VL,
Zen-3D \cite{zen3d}, Zen-Foley \cite{zenfoley}, Zen-Musician \cite{zenmusician}, Zen-Video
\cite{zenvideo}, and the Jin diffusion-transformer family \cite{jin}. Zen5 enters the existing
hanzoai/zen \emph{Spatially-Aware Dynamic Routing} architecture \cite{hanzoai-zen} as the
language-and-reasoning expert, while routing input across modalities to the appropriate sibling,
colocated inside one \texttt{hanzo-engine} binary (Section~\ref{sec:routing}).

\paragraph{MoDE.} As a research direction we introduce \emph{Mixture of Diverse Experts} (MoDE),
a training-time diversity regularizer for Jin's diffusion-transformer MoE \cite{jin}. MoDE produces
experts specialized along the denoising trajectory and across modality-specific latent
subspaces. Section~\ref{sec:mode} describes the objective and an ablation protocol.

\paragraph{Contributions.} The contributions of this work are:

\begin{enumerate}
  \item A practical engine fork (\texttt{zen5}, MIT) that brings DeepSeek V4 Flash to personal Apple
        Silicon hardware with a complete OpenAI / Anthropic / Responses-compatible HTTP surface and
        durable on-disk KV cache.
  \item A reproducible three-stage post-training recipe --- abliteration, identity SFT, agentic SFT ---
        whose ingredients we release alongside the model.
  \item Integration of Zen5 as a routed expert in the existing hanzoai/zen multimodal architecture,
        with parallel sibling-model loading inside a single \texttt{hanzo-engine} process.
  \item MoDE, a diversity-regularized diffusion MoE for unified multimodal generation, with a published
        evaluation protocol against Jin baselines.
\end{enumerate}

\section{Related Work}
\label{sec:related}

\paragraph{DeepSeek V4 Flash and ds4.} DeepSeek V4 Flash \cite{deepseekv4flash} is a 284B-parameter
sparse-MoE language model with multi-token prediction (MTP) draft heads, a heavily compressed indexer
in its attention path, and a 1M-token context ceiling. Sanfilippo's \texttt{ds4} engine
\cite{antirez-ds4} is a self-contained C implementation specialized to this model. It introduces a
two-bit asymmetric routed-experts quantization (\texttt{IQ2\_XXS} up/gate, \texttt{Q2\_K} down) that
leaves shared experts, attention projections, and output untouched, and a SHA1-keyed rendered-prefix
disk KV cache that lets stateless agent clients re-use cached prefixes across process restarts. We
inherit both, and adopt \texttt{ds4}'s Anthropic-compatible \texttt{/v1/messages} endpoint as our
agent-client entry point.

\paragraph{GGML and llama.cpp.} The \texttt{ds4} project is, by Sanfilippo's own attribution
\cite{antirez-ds4}, indebted to \texttt{ggml-org/llama.cpp} \cite{llamacpp} for the GGUF format,
quantization mathematics, and dot-product kernels. Our LICENSE preserves this lineage. We share the
view that local frontier inference is a property of the format $\times$ engine $\times$ agent stack,
and that the GGUF ecosystem is one of the few open standards binding these together.

\paragraph{Refusal directions and abliteration.} Arditi et al.\ \cite{arditi2024refusal} demonstrate
that refusal in instruction-tuned LLMs is mediated by a single linear direction in the residual stream,
and that ablating this direction at inference time produces models that decline to refuse without
otherwise degrading capability. The community has since published \emph{abliteration} pipelines
applying this finding to open-weights models \cite{failspy-abliterate}. \texttt{ds4} itself ships a
\texttt{dir-steering/} subsystem implementing single-vector activation steering at inference time. We
extend this from inference-time steering to a weight-baked abliteration of the routed-MoE-served
upstream, suitable for distribution as GGUF.

\paragraph{Agentic tool calling.} Tool-augmented LLMs have a substantial training-data literature, from
Gorilla \cite{gorilla} and ToolBench \cite{toolbench} to xLAM \cite{xlam} and APIGen \cite{apigen}.
DeepSeek V4 Flash already emits tool calls in DSML \cite{dsml}, an in-band markup language. Our
agentic SFT does not change the protocol; it teaches the model to invoke tools more reliably under
Zen5's particular system prompt and to recover gracefully from common tool-call errors.

\paragraph{Multimodal language models.} The dominant paradigms are late fusion (LLaVA \cite{llava},
LLaVA-NEXT-Interleaved \cite{llavanext}), native multimodal pretraining (Janus
\cite{janus} from DeepSeek itself, Gemini), and modality-expert routing (Flamingo
\cite{flamingo}, MoE-LLaVA \cite{moellava}). The hanzoai/zen architecture \cite{hanzoai-zen} adopts
expert routing: a router selects among heterogeneous specialists for each input. Zen5 enters this
architecture as the language-and-reasoning specialist while preserving the modality experts.

\paragraph{Diffusion-transformer MoE and MoDE.} Diffusion transformers \cite{dit} have become the
dominant architecture for image and video generation; MoE variants \cite{rapid,mome} reduce
inference cost. Jin \cite{jin} combines a diffusion transformer with sparse MoE (2 of 8 experts
active) and a joint embedding space across text, vision, and audio. MoDE in our setting extends Jin
with an explicit \emph{diversity regularizer} over the router assignment matrix, building on the
diversity-of-experts literature in regular MoE \cite{diverse-moe}, with novel application to
denoising-step specialization.

\section{Zen5 Inference Engine}
\label{sec:engine}

\subsection{Architecture}

The Zen5 engine is a fork of \texttt{antirez/ds4} that preserves the upstream sources unchanged.
Operationally, it is a self-contained C and Objective-C executable linked against Apple's Metal
framework on macOS, or against CUDA on Linux. There is no GGML link at runtime, but
\texttt{LICENSE} preserves the GGML authors' copyright per Section~\ref{sec:related}.

Three binaries result from \texttt{make} on macOS:

\begin{itemize}
  \item \texttt{zen5} --- interactive CLI and one-shot prompt mode (\texttt{linenoise}-backed).
  \item \texttt{zen5-server} --- the HTTP server (OpenAI / Anthropic / Responses-compatible) with
        SSE streaming and an on-disk KV cache.
  \item \texttt{zen5-bench} --- a benchmarking harness that measures instantaneous prefill and
        decode rates at fixed context frontiers.
\end{itemize}

These names are symlinks in the working tree to the upstream binary names (\texttt{ds4}, \texttt{ds4-server},
\texttt{ds4-bench}) so that upstream pulls apply cleanly. The Zen5 brand affects packaging, not source.

\subsection{Quantization}

\textbf{The Zen5 weights are the DeepSeek V4 Flash weights, quantized.} We do not train or distill a
distinct model; ``Zen5'' at the weight level is the upstream DeepSeek V4 Flash checkpoint
\cite{deepseekv4flash} converted to GGUF and quantized with the scheme below, optionally with the
behavioral post-training of Section~\ref{sec:posttrain} applied as a small weight delta. There is no
independent zen-5 pretraining run and no separate zen-5 weight set.

DeepSeek V4 Flash at \texttt{IQ2\_XXS}/\texttt{Q2\_K} occupies approximately 81 GB on disk, while
preserving practical agent and tool-calling quality. The asymmetric quantization scheme is critical:
only the routed MoE experts are aggressively quantized (\texttt{IQ2\_XXS} on up/gate matrices,
\texttt{Q2\_K} on down matrices), while shared experts, attention projections, the routing logits,
and the LM head are kept at higher precision. Because routed experts constitute the majority of
model parameters in DeepSeek V4 Flash, this scheme delivers the storage benefit of 2-bit weights
while protecting the components most sensitive to quantization error.

We additionally support an \texttt{imatrix} variant in which weight importance is calibrated against a
held-out corpus. The imatrix variant is the recommended default for Zen5.

\subsection{Disk KV Cache and Stateless-Client Resumption}

A modern coding-agent client (Claude Code, Codex CLI, opencode) re-sends most or all of the prior
conversation on each turn. Without KV reuse this re-prefills tens of thousands of tokens every turn,
even for short edits. The Zen5 engine inherits ds4's KV cache design with three properties:

\begin{enumerate}
  \item \textbf{Rendered-prefix keying.} The cache key is \texttt{SHA1} of the tokenizer-decoded
        \emph{text} of the cached prefix, not the token ID sequence. This lets a request that
        re-tokenizes a generated suffix into a different token count still hit the cache.
  \item \textbf{Save points.} Saves are taken at four moments: \emph{cold} (after first stable
        prefix, before generation), \emph{continued} (at aligned token frontiers during long
        generations), \emph{evict} (before an unrelated session replaces the live checkpoint), and
        \emph{shutdown}.
  \item \textbf{Tool-ID exact replay.} Each tool call gets an unguessable ID; the server stores the
        exact sampled DSML bytes for that ID in a bounded map, persisted in the cache file as a
        \texttt{KTM} tail section. On the next turn, the renderer uses the exact original bytes
        rather than re-rendering from client-side JSON, preserving the byte-exact prefix that the KV
        checkpoint requires.
\end{enumerate}

This subsystem is unmodified from upstream and is the primary reason Zen5 is usable inside multi-turn
agent loops with 100k+ context.

\subsection{HTTP API and SSE Streaming}

\texttt{zen5-server} exposes four wire APIs (Table~\ref{tab:endpoints}). The Anthropic-compatible
\texttt{/v1/messages} endpoint is the recommended entry point for Claude-Code-style clients; in
\texttt{thinking} mode reasoning is streamed in the native API shape rather than mixed into final
text. Tool calls in OpenAI chat streaming and Anthropic messages streaming are emitted as soon as the
DSML invocation is recognized, with parameter bytes forwarded as deltas while generation continues.

\begin{table}[H]
\centering
\begin{tabular}{lll}
\toprule
Wire API & Path & Primary client \\
\midrule
OpenAI Chat Completions & \texttt{POST /v1/chat/completions} & opencode, Pi \\
OpenAI Completions & \texttt{POST /v1/completions} & legacy \\
OpenAI Responses & \texttt{POST /v1/responses} & Codex CLI \\
Anthropic Messages & \texttt{POST /v1/messages} & Claude Code \\
\bottomrule
\end{tabular}
\caption{Zen5 HTTP wire APIs.}
\label{tab:endpoints}
\end{table}

\subsection{Performance}

Table~\ref{tab:speed} reproduces the upstream single-run Metal CLI numbers with
\texttt{--ctx 32768}, \texttt{--nothink}, greedy decoding, and \texttt{-n 256}. We measured a smoke
test on M4 Max 128 GB: cold residency 22.6 s, first-token prefill 27.30 tok/s on a 32-token prompt.
Full sweep numbers on M4 Max will be reported with the public release.

\begin{table}[H]
\centering
\begin{tabular}{llrrr}
\toprule
Machine & Quant & Prompt & Prefill & Generation \\
\midrule
MacBook Pro M3 Max, 128 GB & q2-imatrix & short & 58.52 t/s & 26.68 t/s \\
MacBook Pro M3 Max, 128 GB & q2-imatrix & 11709 tok & 250.11 t/s & 21.47 t/s \\
Mac Studio M3 Ultra, 512 GB & q2-imatrix & short & 84.43 t/s & 36.86 t/s \\
Mac Studio M3 Ultra, 512 GB & q4-imatrix & 12018 tok & 448.82 t/s & 26.62 t/s \\
DGX Spark GB10, 128 GB & q2-imatrix & 7047 tok & 343.81 t/s & 13.75 t/s \\
\bottomrule
\end{tabular}
\caption{Reproduced from upstream \texttt{antirez/ds4}; M4 Max numbers in progress.}
\label{tab:speed}
\end{table}

\subsection{Planned Engine Extensions}

While the upstream engine is functionally complete for single-model inference, the Zen5 multimodal
roadmap (Section~\ref{sec:routing}) requires several engine-level capabilities not yet present
upstream. We track these in \texttt{PUNCH\_LIST\_FROM\_DS4.md} in our companion \texttt{hanzo-engine}
crate; the highest-priority items are durable disk KV across non-DeepSeek models, a generic
modality router transport, and concurrent multi-model loading with VRAM budgeting.

\section{Post-Training: Identity, Alignment, and Agentic Tool Use}
\label{sec:posttrain}

The base DeepSeek V4 Flash checkpoint requires three post-training stages before it can be released
as Zen5. We apply them in the order \emph{abliteration $\rightarrow$ identity SFT $\rightarrow$
agentic SFT}, because abliteration operates on the base persona's refusal substrate and is most clean
before identity is rewritten.

\subsection{Stage 1: Abliteration}

\paragraph{Goal.} Reduce unhelpful refusals on legitimate dual-use prompts (security analysis, code
review of obfuscated samples, protocol-level explanations) without removing the model's ability to
recognize and decline genuinely harmful requests.

\paragraph{Method.} Following Arditi et al.\ \cite{arditi2024refusal}, we identify a one-dimensional
subspace $\mathbf{r}$ in the residual stream that mediates refusal:

\begin{equation}
\mathbf{r}^{(\ell)} = \frac{1}{|D_\text{harm}|} \sum_{x \in D_\text{harm}} h^{(\ell)}_x
- \frac{1}{|D_\text{benign}|} \sum_{x \in D_\text{benign}} h^{(\ell)}_x
\end{equation}

where $h^{(\ell)}_x$ is the post-layer-$\ell$ hidden state at the final input token for prompt $x$,
$D_\text{harm}$ is a 64-prompt set of refusal-eliciting requests, and $D_\text{benign}$ is a
matched-distribution non-refusing set. We select the layer $\ell^*$ that maximizes the cosine
distance between $\mathbf{r}$ projections of harmful vs.\ benign held-out prompts. We then orthogonalize
the down-projection matrices of all routed and shared expert MLPs against $\mathbf{r}^{(\ell^*)}$:

\begin{equation}
W_\text{down}^{\text{post}} \;=\; W_\text{down} - W_\text{down} \frac{\mathbf{r}\mathbf{r}^\top}{\|\mathbf{r}\|^2}
\end{equation}

The orthogonalization is weight-baked (no inference-time hooks), so the resulting GGUF is a drop-in
replacement.

\paragraph{Evaluation.} We measure (a) refusal rate on a 200-prompt evaluation split spanning the
Cybersecurity, Biorisk, and General Dual-Use categories, and (b) helpfulness retention on MT-Bench
and IFEval. Target: refusal rate $\leq$ 15\% on dual-use without MT-Bench drop greater than 2 points.
Inference-time steering via the engine's \texttt{dir-steering/} subsystem remains available as a
runtime override.

\paragraph{Safety considerations.} We do not abliterate the model on prompts touching CSAM,
mass-casualty CBRN synthesis routes, or specific real-person targeting. The abliteration corpus
$D_\text{harm}$ is filtered to exclude these categories so that the refusal direction we ablate is
the over-refusal direction, not the entire safety substrate. Post-abliteration we re-evaluate against
WMDP \cite{wmdp} and our internal red-team suite; the model that ships under the Zen5 brand must
preserve the original refusal rate on these categories. This boundary is enforced by gating the
release on red-team pass-fail, not by trust in the procedure alone.

\subsection{Stage 2: Identity Supervised Fine-Tuning}

\paragraph{Goal.} Replace the model's self-description from ``DeepSeek V4 Flash'' to ``Zen5'' across
identity-eliciting prompts (\emph{Who are you?}, \emph{What model is this?}, \emph{What is your
training cutoff?}, \emph{What can you do?}), without altering capability.

\paragraph{Method.} A small ($\sim$10k examples) supervised fine-tuning corpus
generated from the Zen brand guide and the Zen family overview \cite{zenfamily}.
Training: LoRA \cite{lora} at rank 32 on the routed-experts up/gate/down projections, 3 epochs,
learning rate $3 \times 10^{-5}$, cosine schedule. LoRA is the right tool here because identity is a
shallow rewrite; full fine-tuning is unnecessary and harms downstream tool-call quality.

\paragraph{Evaluation.} 100 held-out identity probes; required pass rate $\geq$ 95\%. MT-Bench /
IFEval drop tolerance $\leq$ 1 point.

\subsection{Stage 3: Agentic Tool-Calling SFT --- the Zen Agentic Dataset}

\paragraph{Goal.} Increase tool-call success rate, decrease tool-call format errors, and improve
recovery behavior when a tool returns an error or empty result.

\paragraph{Dataset.} The Zen Agentic Dataset comprises:

\begin{itemize}
  \item $\sim$50k single-tool single-turn examples seeded from xLAM \cite{xlam} and APIGen
        \cite{apigen}, rewritten into DSML.
  \item $\sim$10k multi-tool multi-turn agent traces from successful runs of our internal Zen-Coder
        and Zen-Agent harnesses on real coding tasks, filtered by unit-test pass.
  \item $\sim$5k \emph{recovery} examples: tool errors, malformed JSON returns, ambiguous results, in
        which the assistant must reason about the failure and either retry or escalate to the user.
\end{itemize}

The dataset is released under CC-BY-4.0 with provenance for each example.

\paragraph{Training.} Two-stage: first an SFT pass over the full mixture for 2 epochs, then a
DPO \cite{dpo} pass on a 20k-pair preference set built by sampling two completions per prompt and
labelling them by tool-call success.

\paragraph{Evaluation.} ToolBench \cite{toolbench} task-pass rate, SWE-Bench-Lite agent pass rate, and
an internal Zen-Coder regression on 500 held-out coding tasks. Targets: $\geq$ 70\% ToolBench,
$\geq$ 35\% SWE-Bench-Lite agentic, with no drop on MT-Bench.

\section{Multimodal Routing and Sibling Experts}
\label{sec:routing}

\subsection{Architectural Position}

The hanzoai/zen \emph{Spatially-Aware Dynamic Routing} architecture \cite{hanzoai-zen} defines a router
that connects specialized expert models for vision, language, codegen, multimodal reasoning, and 3D
spatial inputs. In the integration described here, Zen5 (i.e.\ the DeepSeek V4 Flash repackage) is
intended to serve as the \emph{language-and-reasoning expert}. The sibling experts named below are
documented in their own reports; several are themselves built on upstream open models rather than
trained from scratch, and the routing integration is a design described here rather than a shipped,
benchmarked system (Section~\ref{sec:experiments} marks routing evaluation as pending).

\begin{figure}[H]
\centering
\begin{tabular}{|c|c|c|}
\hline
\textbf{Modality} & \textbf{Sibling expert} & \textbf{Backend} \\
\hline
Language, reasoning, tools & Zen5 (DeepSeek V4 Flash, abliterated) & \texttt{zen5-server} \\
Vision-language & Zen-VL 4B/8B/30B & mistral-rs (vision adapter) \\
3D understanding & Zen-3D & native PyTorch / mistral-rs \\
Audio recognition & Whisper-derived & mistralrs-audio \\
Foley / sound design & Zen-Foley & native PyTorch \\
Music generation & Zen-Musician & native PyTorch \\
Video generation & Zen-Video & diffusion (Jin lineage) \\
Image / diffusion gen & Jin / Zen-Artist & Jin (diffusion-MoE) \\
\hline
\end{tabular}
\caption{Sibling expert roster. Several siblings are themselves implemented on top of \texttt{mistral-rs};
the Zen5 reasoning core uses the \texttt{zen5-server} runtime.}
\end{figure}

\subsection{Modality Router}

The router operates on the request boundary and decides, per request, which sibling handles which
content. Two versions:

\paragraph{V1 --- rule-based.} For the public release of Zen5 we ship a deterministic router that
inspects message content types (\texttt{image\_url}, \texttt{audio}, \texttt{video}, \texttt{model3d})
in the incoming OpenAI / Anthropic request, plus a small set of regex- and prefix-based intent rules
on the text (e.g., a request beginning with ``Render an image of\ldots'' routes to Zen-Artist; a
request containing a \texttt{.obj} or \texttt{.glb} attachment routes to Zen-3D). The Zen5 reasoning
core acts as orchestrator: it sees a summary of the sibling's response and integrates it into the
final reply. This corresponds to the \emph{tool-call} pattern but with siblings rather than external
APIs.

\paragraph{V2 --- learned.} A two-tower small classifier (\textasciitilde 50M parameters) trained on
routing traces collected from V1. Input: the full conversation prefix plus attachment metadata.
Output: a multi-label sibling assignment plus a confidence. The classifier is trained with the
standard cross-entropy objective on V1-labelled examples and a held-out human-validated set. We
adopt the \emph{learned-routing-policy} component of the hanzoai/zen \texttt{RoutingPolicy}
\cite{hanzoai-zen} as the training target.

\subsection{Parallel Multi-Model Loading in \texttt{hanzo-engine}}

The current \texttt{mistralrs-server-core} runtime supports model lifecycle endpoints
(\texttt{/v1/models/\{load,unload,reload,tune\}}) but serializes inference through a single live
graph. For Zen5's modality routing to work without per-request reloads, the engine must hold multiple
expert models simultaneously and dispatch concurrent requests. Our planned extension introduces a
\texttt{ModelRegistry} layer above the existing engine handle, with three properties:

\begin{enumerate}
  \item \textbf{VRAM budgeting.} Each registered model declares its peak resident size at the
        chosen quant. The registry refuses loads that would exceed the wired GPU memory limit
        (\texttt{iogpu.wired\_limit\_mb} on macOS). Default Zen5 deployment: Zen5 at 81 GB plus a
        Zen-VL-8B at 5 GB plus Jin diffusion-MoE at 8 GB leaves the user 26 GB headroom on a 120 GB
        wired limit.
  \item \textbf{Per-model async queues.} Each model has its own request queue and graph worker;
        cross-model parallelism is bound only by GPU compute and memory bandwidth.
  \item \textbf{Shared KV semantics.} Disk KV cache is keyed by \texttt{(model\_id, sha1(rendered\_text))}
        so that a sibling switch on the next turn does not cross-corrupt cached prefixes.
\end{enumerate}

The \texttt{zen5-server} runtime stays single-model; multi-model multiplexing lives in
\texttt{hanzo-engine}, which already wraps \texttt{mistral-rs} and is the natural location for the
sibling roster.

\subsection{Zoo Desktop Integration}

The user-facing surface is Zoo Desktop \cite{zoodesktop}, a Tauri 2.x application that already
supports adding LLM providers with a custom \texttt{external\_url} per provider. The Zen5 + siblings
deployment is exposed to Zoo Desktop as a single provider entry pointing at
\texttt{http://127.0.0.1:36900/v1}; routing across siblings is invisible to the desktop client. From
the Zoo Desktop user's perspective there is one Zen model that handles text, vision, audio, video,
and 3D --- the modality split is an engine-internal optimization.

\section{MoDE: Mixture of Diverse Experts for Diffusion Generation}
\label{sec:mode}

\subsection{Motivation}

The Jin diffusion-transformer MoE \cite{jin} activates 2 of 8 experts per token under a load-balanced
top-$k$ gating function. The standard load-balancing loss \cite{switch} encourages roughly equal
expert utilization but does not constrain experts to specialize in functionally distinct directions
--- experts can converge to near-identical computations under the load-balance prior alone. In
diffusion generation, where the model must perform substantially different work at $t=999$ (early,
coarse-structure denoising) and $t=10$ (late, fine-detail refinement), and across modality-distinct
latent regions (audio spectrogram patches versus image patches versus point-cloud patches), forcing
\emph{diverse} expert specialization is intuitively valuable.

\subsection{Objective}

Let $G \in \mathbb{R}^{N \times E}$ be the gating assignment matrix over a batch of $N$ tokens and
$E$ experts, where $G_{ne} = 1$ if token $n$ is routed to expert $e$. The MoDE training objective
augments the Jin loss with a diversity term:

\begin{equation}
\mathcal{L}_\text{MoDE} \;=\; \mathcal{L}_\text{denoise} \;+\; \lambda_\text{lb}\, \mathcal{L}_\text{lb} \;+\;
\lambda_\text{div}\, \mathcal{L}_\text{div}
\end{equation}

\noindent where $\mathcal{L}_\text{denoise}$ is the standard diffusion noise-prediction loss,
$\mathcal{L}_\text{lb}$ is the load-balancing loss \cite{switch}, and the diversity term
$\mathcal{L}_\text{div}$ is the negative pairwise expert-direction Frobenius distance:

\begin{equation}
\mathcal{L}_\text{div} \;=\; -\frac{2}{E(E-1)} \sum_{i < j} \big\| W^{(i)} - W^{(j)} \big\|_F^2
\end{equation}

\noindent where $W^{(i)}$ is the concatenated MLP weight matrix of expert $i$. This is a simple,
weight-space diversity prior; it encourages experts to occupy distinct directions in parameter space.
Alternative forms include input-conditioned diversity (computing expert-output disagreement across
the batch) and a denoising-step conditional diversity (encouraging different gating at different
$t$). We will report ablations across these forms.

\subsection{Specialization Hypotheses}

We expect three forms of specialization to emerge from $\mathcal{L}_\text{div}$:

\begin{enumerate}
  \item \textbf{Denoising-step specialization.} Some experts will preferentially activate on early
        diffusion steps (coarse structure) and others on late (fine detail).
  \item \textbf{Modality specialization.} In Jin's joint embedding space the same expert pool serves
        text, vision, and audio tokens; we expect diversity-regularized experts to partition by
        modality.
  \item \textbf{Latent-region specialization.} Within a single modality, different latent regions
        (e.g., faces vs.\ backgrounds in image patches) will route to different experts.
\end{enumerate}

We will measure these by post-hoc clustering of the gating assignment matrix and visualizing per-expert
activation maps.

\subsection{Evaluation Protocol}

We compare MoDE against the standard Jin baseline on:

\begin{itemize}
  \item Image generation: FID on COCO-30K, FID on ImageNet-256.
  \item Audio generation: FAD on AudioSet, prompt-conditioned MOS on a 200-sample human eval.
  \item 3D generation: CLIP-similarity for text-to-3D on Cap3D.
  \item Multimodal mixing: a held-out set of cross-modal prompts (image-conditioned audio,
        audio-conditioned image, etc.) judged by human annotators.
\end{itemize}

The diversity-regularization sweep covers $\lambda_\text{div} \in \{0, 10^{-3}, 10^{-2}, 10^{-1}\}$
and one ablation removing $\mathcal{L}_\text{lb}$ entirely (to test whether diversity subsumes
load-balancing). Compute budget: training MoDE-Jin-Small (1.5B active, 8 experts of 220M) for 100k
steps on 16 H100s.

\section{Experiments and Evaluation}
\label{sec:experiments}

This section will be populated as the public Zen5 release lands; the framework is described here so
that third-party replications can begin in parallel with our own.

\paragraph{Engine smoke test.} On M4 Max 128 GB, Apple Metal, \texttt{q2-imatrix},
\texttt{ctx=32768}, \texttt{--nothink}: cold residency 22.6 s, prefill 27.30 tok/s on a 32-token
prompt, single-token decode latency 112 ms. Full speed sweeps to be reported alongside Table~\ref{tab:speed}.

\paragraph{Quality.} Pending. We will report MT-Bench, IFEval, MMLU, HumanEval, GPQA, SWE-Bench-Lite,
WMDP-Cyber, and our internal Zen-Coder 500-task agentic regression at the public release point. We
will additionally report deltas pre- vs.\ post-abliteration on the same suite to demonstrate
helpfulness retention.

\paragraph{Routing.} The V1 rule-based router will be evaluated on a 1000-example multimodal request
set with human-labelled ground-truth sibling assignment. The V2 learned router will be evaluated on
the same set with a target $\geq$ 95\% top-1 sibling accuracy.

\paragraph{MoDE.} Pending. We will report FID / FAD / CLIP-similarity vs.\ Jin baseline and ablations
across $\lambda_\text{div}$.

\section{Future Work}
\label{sec:future}

\paragraph{Beyond q2-imatrix.} The current quant is the floor for single-laptop deployment. We expect
DeepSeek to release updated V4 Flash checkpoints; we will track them. For Mac Studio M3 Ultra 512 GB
the q4-imatrix variant offers the same context with substantially higher quality and is the better
default for that hardware class.

\paragraph{Distributed inference.} The engine is single-host today. A future direction is sharding
the expert pool across two M-series machines connected by Thunderbolt 5, with the indexer and shared
experts on one host and routed-experts split across both. This is plausible because routed-experts
exhibit token-level sparsity that maps well to inter-host traffic.

\paragraph{Continuous routing learning.} The V2 learned router can be trained continuously on
production routing decisions. We will release a privacy-respecting export protocol so that operators
can opt into contributing routing traces.

\paragraph{MoDE at scale.} A 30B-active MoDE-Jin variant is the natural follow-on to the 1.5B
ablation \emph{proposed} in Section~\ref{sec:mode}. That ablation has not been run; hardware
permitting, we will report it (and any scale-up) in a follow-up whitepaper.

\section{Conclusion}

Zen5 is not a new foundation model. It is a re-packaging, a behavioral re-alignment, and a
multimodal routing layer over a foundation that was already public --- DeepSeek V4 Flash --- using an
engine that was already public --- \texttt{antirez/ds4}. The contribution is in the integration: a
single \texttt{hanzo-engine} process that holds a 284B reasoning model alongside vision, 3D, audio,
and diffusion siblings; a behavioral pipeline that converts a generic upstream into a usable Zen5
persona without retraining; and a research direction (MoDE) that aims to unify multimodal generation
under one diffusion-MoE backbone. We ship as alpha, in public, with attribution intact.

\section*{Acknowledgements}

This work would not exist without \textbf{Salvatore Sanfilippo} (antirez) and the contributors to
\texttt{ds4}, and without the \texttt{ggml-org/llama.cpp} community on whose engineering it stands.
Both copyrights are preserved unchanged in our LICENSE. We additionally thank the DeepSeek team for
releasing V4 Flash with permissive terms, and the authors of the refusal-direction paper for the
foundation that makes Stage 1 of Section~\ref{sec:posttrain} possible.

\begin{thebibliography}{99}

\bibitem{deepseekv4flash}
DeepSeek-AI.
\newblock {\em DeepSeek V4 Flash Technical Report}, 2026.

\bibitem{antirez-ds4}
S.~Sanfilippo.
\newblock \texttt{antirez/ds4}: A self-contained inference engine for DeepSeek V4 Flash, 2026.
\newblock \url{https://github.com/antirez/ds4}.

\bibitem{llamacpp}
G.~Gerganov et al.
\newblock \texttt{ggml-org/llama.cpp}: Inference of Llama-family models in pure C/C++, 2023--2026.
\newblock \url{https://github.com/ggml-org/llama.cpp}.

\bibitem{arditi2024refusal}
A.~Arditi, O.~Obeso, A.~Sylvestre, K.~Bian, A.~Conmy, N.~Nanda.
\newblock Refusal in Language Models Is Mediated by a Single Direction.
\newblock {\em arXiv preprint arXiv:2406.11717}, 2024.

\bibitem{failspy-abliterate}
B.~Schmidt et al.
\newblock \texttt{failspy/abliterator}: Practical abliteration of open-weight language models, 2024.

\bibitem{gorilla}
S.~Patil, T.~Zhang, X.~Wang, J.~E.~Gonzalez.
\newblock Gorilla: Large Language Model Connected with Massive APIs.
\newblock {\em NeurIPS}, 2023.

\bibitem{toolbench}
Y.~Qin et al.
\newblock ToolLLM: Facilitating Large Language Models to Master 16000+ Real-world APIs.
\newblock {\em ICLR}, 2024.

\bibitem{xlam}
J.~Zhang et al.
\newblock xLAM: A Family of Large Action Models to Empower AI Agent Systems.
\newblock {\em arXiv preprint arXiv:2409.03215}, 2024.

\bibitem{apigen}
Z.~Liu et al.
\newblock APIGen: Automated Pipeline for Generating Verifiable and Diverse Function-Calling Datasets.
\newblock {\em NeurIPS}, 2024.

\bibitem{dsml}
DeepSeek-AI.
\newblock DSML: The DeepSeek Markup Language for tool invocations, 2026.

\bibitem{llava}
H.~Liu, C.~Li, Q.~Wu, Y.~J.~Lee.
\newblock Visual Instruction Tuning.
\newblock {\em NeurIPS}, 2023.

\bibitem{llavanext}
F.~Li, H.~Zhang, P.~Zhang, S.~Zhang, T.~Lu.
\newblock LLaVA-NeXT-Interleaved.
\newblock 2024.

\bibitem{janus}
W.~Xu et al.
\newblock Janus: Decoupling Visual Encoding for Unified Multimodal Understanding and Generation.
\newblock DeepSeek-AI, 2024.

\bibitem{flamingo}
J.-B.~Alayrac et al.
\newblock Flamingo: a Visual Language Model for Few-Shot Learning.
\newblock {\em NeurIPS}, 2022.

\bibitem{moellava}
B.~Lin et al.
\newblock MoE-LLaVA: Mixture of Experts for Large Vision-Language Models.
\newblock {\em arXiv preprint arXiv:2401.15947}, 2024.

\bibitem{dit}
W.~Peebles, S.~Xie.
\newblock Scalable Diffusion Models with Transformers.
\newblock {\em ICCV}, 2023.

\bibitem{rapid}
Z.~Sun et al.
\newblock Rapid-MoE Diffusion Transformer, 2024.

\bibitem{mome}
W.~Wang et al.
\newblock Image as a Foreign Language: BEiT-3 and Mixture-of-Modality-Experts.
\newblock {\em CVPR}, 2023.

\bibitem{diverse-moe}
A.~Ruiz, S.~Sun et al.
\newblock Diversity-Regularized Mixture-of-Experts, 2023.

\bibitem{switch}
W.~Fedus, B.~Zoph, N.~Shazeer.
\newblock Switch Transformer: Scaling to Trillion Parameter Models with Simple and Efficient Sparsity.
\newblock {\em JMLR}, 2022.

\bibitem{dpo}
R.~Rafailov et al.
\newblock Direct Preference Optimization: Your Language Model is Secretly a Reward Model.
\newblock {\em NeurIPS}, 2023.

\bibitem{lora}
E.~J.~Hu et al.
\newblock LoRA: Low-Rank Adaptation of Large Language Models.
\newblock {\em ICLR}, 2022.

\bibitem{wmdp}
N.~Li et al.
\newblock The WMDP Benchmark: Measuring and Reducing Malicious Use With Unlearning.
\newblock 2024.

\bibitem{zen3d}
Zen Team.
\newblock Zen-3D: Spatially-Aware Generation and Understanding.
\newblock {\em Zen Whitepapers}, 2025.

\bibitem{zenfoley}
Zen Team.
\newblock Zen-Foley: Multimodal Sound Design.
\newblock {\em Zen Whitepapers}, 2025.

\bibitem{zenmusician}
Zen Team.
\newblock Zen-Musician: Music Generation in the Zen Family.
\newblock {\em Zen Whitepapers}, 2025.

\bibitem{zenvideo}
Zen Team.
\newblock Zen-Video: Text-to-Video and Image-to-Video Generation.
\newblock {\em Zen Whitepapers}, 2025.

\bibitem{jin}
Hanzo AI / zenlm.
\newblock Jin: Diffusion-Transformer Mixture-of-Experts with Joint Embedding Space.
\newblock \url{https://github.com/hanzoai/jin}, 2026.

\bibitem{hanzoai-zen}
Hanzo AI.
\newblock Zen: Spatially-Aware Dynamic Routing Architecture.
\newblock Internal technical document, 2025.

\bibitem{zenfamily}
Zen Team.
\newblock Zen AI Model Family Overview.
\newblock {\em Zen Whitepapers}, 2026.

\bibitem{zoodesktop}
Zoo Labs Foundation.
\newblock Zoo Desktop: Multimodal AI Workstation.
\newblock \url{https://github.com/zooai/app}, 2026.

\end{thebibliography}

\end{document}
